\section{Reset Model}

Reset places the processor into a known architectural state and begins instruction execution
from the reset vector.

VIV-32 v\version{} defines only machine-mode reset. User-mode reset, supervisor-mode
reset, secure boot, and multi-stage privilege transitions are reserved for future versions.

\section{Reset Vector}

On reset, the program counter is set to the architectural reset vector.

\begin{center}
\code{\%pc = 0x00000000}
\end{center}

\noindent
The reset vector address must be 4-byte aligned. The platform must ensure that
executable reset code is available at the reset vector address.

\noindent
A platform may implement the reset vector using ROM, RAM preloaded by an
emulator, memory-mapped boot firmware, or another platform-defined boot mechanism.

\section{Reset State}

After reset, architectural state is initialized as follows:

\begin{center}
\begin{tabular}{lll}
\hline
State & Reset value & Notes \\
\hline
\code{PC} & \code{0x00000000} & Reset vector \\
\code{\$0} & \code{0x00000000} & Hardwired zero \\
\code{\$1--\$15} & Undefined & Software must initialize before use \\
\code{\#0--\#7} & Undefined & Software must initialize before use \\
\code{\%sr.N} & \code{0} & Negative flag cleared \\
\code{\%sr.Z} & \code{0} & Zero flag cleared \\
\code{\%sr.C} & \code{0} & Carry flag cleared \\
\code{\%sr.V} & \code{0} & Overflow flag cleared \\
\code{\%sr.I} & \code{0} & Infinity flag cleared \\
\code{\%sr.U} & \code{0} & Underflow flag cleared \\
\code{\%sr.X} & \code{0} & Inexact flag cleared \\
\code{\%sr.E} & \code{0} & Arithmetic error flag cleared \\
\code{\%sr.IE} & \code{0} & Interrupts disabled \\
\code{\%epc} & Undefined & No exception has been taken \\
\code{\%esr} & Undefined & No exception status has been saved \\
\code{\%ecause} & \code{0x00} & Reset cause \\
\code{\%edata} & \code{0x00000000} & No exception data \\
\code{\%evbase} & Platform-defined & Must be 16-byte aligned; software must initialize
 before exception handling is used \\
\hline
\end{tabular}
\end{center}

\noindent
Any platform-defined reset value of \code{\%evbase} must satisfy the architectural 16-byte
alignment requirement.
General-purpose registers other than \code{r0} are undefined after reset. Boot code
must initialize \code{\$sp}, global pointers, frame pointers, and any other registers
required by the ABI or execution environment.

\section{Reset Cause}

Reset is represented architecturally by exception cause value \code{0x00}.
After reset, \code{\%ecause} contains \code{0x00}.

\noindent
Reset does not require the processor to enter through the ordinary exception
vector calculation. On reset, execution begins directly at the reset vector.

\section{Initial Interrupt State}

Interrupts are disabled after reset.

\begin{itemize}
\item The \code{IE} bit in \code{sr} is cleared.
\item Timer interrupts are not taken until software enables interrupts.
\item External interrupts are not taken until software enables interrupts.
\item Pending interrupt state, if any, is platform-defined.
\end{itemize}

\noindent
Boot firmware should initialize \code{evbase}, configure platform interrupt devices, install exception handlers, and then enable interrupts using \code{ei} when ready.

\section{Initial Memory State}

The contents of RAM after reset are platform-defined.

\begin{itemize}
\item RAM may contain undefined data.
\item ROM or boot firmware contents are platform-defined.
\item Memory-mapped device reset state is defined by the platform specification.
\item The base ISA does not require memory to be cleared on reset.
\end{itemize}

\noindent
Software must not assume that RAM is zero-filled unless the platform specification explicitly guarantees it.

\section{Boot Responsibilities}

The base architecture defines only the architectural reset state. Platform firmware or emulator boot code is responsible for preparing the execution environment.

Typical boot responsibilities include:

\begin{itemize}
\item initializing the stack pointer;
\item initializing \code{evbase};
\item installing exception and interrupt handlers;
\item clearing or initializing the \code{.bss} region;
\item initializing global data;
\item configuring memory-mapped devices;
\item configuring timer and external interrupt sources;
\item enabling interrupts when the system is ready;
\item transferring control to firmware, a monitor, a kernel, or a standalone program.
\end{itemize}

\section{Boot Sequence Example}

A simple platform boot sequence may perform the following operations:

\begin{enumerate}
\item Begin execution at \code{0x00000000}.
\item Disable interrupts using \code{di}.
\item Initialize \code{sp}.
\item Initialize \code{evbase}.
\item Install or verify exception vector contents.
\item Initialize memory and global data.
\item Configure platform devices.
\item Enable interrupts using \code{ei}, if required.
\item Transfer control to the main program.
\end{enumerate}

\noindent
This sequence is illustrative. The exact boot sequence is defined by the platform and execution environment.
